\chapter{Formulace problému}

\section{Účel}

Cílem tohoto projektu je teoreticky navrhnout informační systém pro správu skladových
součástek. Systém má umožňovat evidenci součástek, jejich kategorií, dodavatelů
a skladových lokací. Dále má podporovat sledování pohybů součástek ve skladu,
jako je naskladnění, vyskladnění nebo přesun mezi jednotlivými lokacemi.

Součástí systému je také možnost přiřazovat součástky k projektům a evidovat
jejich množství potřebné pro daný projekt.

\section{Specifikace}

Návrh systému je znázorněn pomocí UML diagramů, které popisují jednotlivé třídy,
jejich atributy, vztahy mezi nimi a chování systému při klíčových operacích.
Systém zahrnuje objektový model (class diagram), dynamický model (stavové diagramy)
a model interakcí (use case diagram, scénáře, sekvenční a aktivity diagram).

\chapter{Datový slovník}

\section{Soucastka}

Skladová položka evidovaná v~systému WMS s~aktuálním množstvím na skladě. Atribut rezervovaneMnozstvi sleduje, kolik kusů je aktuálně blokováno pro projekty -- součástka tak může být Dostupna i s~částečnou rezervací zároveň. Součástka může nabývat stavů: Dostupna (k~dispozici, případně s~rezervovanou částí), Vydana (fyzicky vydána ze skladu) nebo Odepsana (trvale vyřazena z~evidence).

\textbf{Atributy:}\begin{itemize}[noitemsep,topsep=2pt]\item id: int\item nazev: String\item mnozstvi: int\item rezervovaneMnozstvi: int\item stav: String (Dostupna $|$ Vydana $|$ Odepsana)\end{itemize}

\textbf{Metody:}\begin{itemize}[noitemsep,topsep=2pt]\item zvysMnozstvi(pocet: int)\item snizMnozstvi(pocet: int)\item rezervovat(pocet: int)\item zrusitRezervaci(pocet: int)\item overDostupnost(): boolean\end{itemize}

\section{Projekt}

Projekt, ke kterému jsou přiřazeny součástky pro sledování spotřeby.

\textbf{Atributy:}\begin{itemize}[noitemsep,topsep=2pt]\item id: int\item nazev: String\end{itemize}

\textbf{Metody:}\begin{itemize}[noitemsep,topsep=2pt]\item priradSoucastku(Soucastka)\item odeberSoucastku(Soucastka)\end{itemize}

\section{SkladovyPohyb}

Abstraktní třída reprezentující pohyb součástek -- příjem, výdej nebo přesun. Každý pohyb prochází životním cyklem: vytvoření, validace dat a provedení pohybu.

\textbf{Atributy:}\begin{itemize}[noitemsep,topsep=2pt]\item id: int\item datum: Date\item stav: String (Dostupna $|$ Vydana $|$ Odepsana)\item typ: String (prijem $|$ vydej $|$ presun)\item mnozstvi: int\end{itemize}

\textbf{Metody:}\begin{itemize}[noitemsep,topsep=2pt]\item \{abstract\} vytvoritPohyb()\item overitData()\item provestPohyb()\end{itemize}

\section{PrijemPohyb}

Konkrétní pohyb typu příjem -- naskladnění součástek od dodavatele. Dědí ze SkladovyPohyb.

\textbf{Atributy:}\begin{itemize}[noitemsep,topsep=2pt]\item dodavatel: String – název dodavatele součástek\item dědí ze SkladovyPohyb: id, datum, stav, typ, mnozstvi\end{itemize}

\textbf{Metody:}\begin{itemize}[noitemsep,topsep=2pt]\item vytvoritPohyb()\end{itemize}

\section{VydejPohyb}

Konkrétní pohyb typu výdej -- vyskladnění součástek ze skladu. Dědí ze SkladovyPohyb. Atribut duvod popisuje účel výdeje, např. použití v~projektu, reklamace, interní spotřeba nebo oprava.

\textbf{Atributy:}\begin{itemize}[noitemsep,topsep=2pt]\item duvod: String – účel výdeje (projekt, reklamace, interní spotřeba, oprava)\item dědí ze SkladovyPohyb: id, datum, stav, typ, mnozstvi\end{itemize}

\textbf{Metody:}\begin{itemize}[noitemsep,topsep=2pt]\item vytvoritPohyb()\end{itemize}

\section{PresunutiPohyb}

Konkrétní pohyb typu přesun -- přemístnění součástek mezi lokacemi ve skladu. Dědí ze SkladovyPohyb.

\textbf{Atributy:}\begin{itemize}[noitemsep,topsep=2pt]\item zdrojLokaceId: int – ID zdrojové lokace\item cilLokaceId: int – ID cílové lokace\item dědí ze SkladovyPohyb: id, datum, stav, typ, mnozstvi\end{itemize}

\textbf{Metody:}\begin{itemize}[noitemsep,topsep=2pt]\item vytvoritPohyb()\item zmenLokaci()\end{itemize}

\section{Sklad}

Fyzický sklad rozdělený na lokace, ve kterých jsou uloženy součástky.

\textbf{Atributy:}\begin{itemize}[noitemsep,topsep=2pt]\item id: int\item nazev: String\end{itemize}

\textbf{Metody:} najdiSoucastku(id: int): Soucastka

\section{Lokace}

Konkrétní místo v~skladu identifikované označením (např. A-01-03).

\textbf{Atributy:}\begin{itemize}[noitemsep,topsep=2pt]\item id: int\item oznaceni: String\end{itemize}

\section{Uzivatel}

Osoba spravující součástky, pohyby, projekty a sklady v~systému.

\textbf{Atributy:}\begin{itemize}[noitemsep,topsep=2pt]\item id: int\item jmeno: String\item email: String\end{itemize}

\newpage
\chapter{Objektový model}

Diagram zachycuje strukturu systému WMS a vztahy mezi jeho hlavními třídami.
Součástky jsou přiřazovány k~projektům a fyzicky umístěny v~lokacích skladu.
Třída Uzivatel reprezentuje uživatelskou roli v~systému. Vytvoření pohybu provádí voláním SkladovyPohyb.vytvoritPohyb() a přiřazení součástky k~projektu voláním Projekt.priradSoucastku().

Diagram využívá \textbf{generalizaci}: abstraktní třída SkladovyPohyb je základem pro tři konkrétní podtřídy
PrijemPohyb (naskladnění), VydejPohyb (vyskladnění) a PresunutiPohyb (přesun mezi lokacemi),
které dědí společné atributy a přidávají vlastní specifické atributy.

Třída SkladovyPohyb deklaruje abstraktní metodu vytvoritPohyb(), kterou každá podtřída implementuje po svém:
\begin{itemize}
  \item PrijemPohyb.vytvoritPohyb() -- vytvoří příjem a nastaví atribut dodavatel.
  \item VydejPohyb.vytvoritPohyb() -- vytvoří výdej a nastaví atribut duvod.
  \item PresunutiPohyb.vytvoritPohyb() -- vytvoří přesun a nastaví zdrojLokaceId a cilLokaceId.
\end{itemize}

Třída SkladovyPohyb zaznamenává každou změnu stavu zásob v~systému. Konkrétně:
\begin{itemize}
  \item PrijemPohyb -- příjem součástek na sklad od dodavatele; zvyšuje množství dané součástky.
  \item VydejPohyb -- výdej součástek ze skladu na projekt; snižuje množství dané součástky.
  \item PresunutiPohyb -- přemístnění součástek mezi lokácemi ve skladu; mění přiřazenou lokaci součástky.
\end{itemize}
Každý pohyb se váže na konkrétní součástku (vazba je součástí, multiplicita $* \rightarrow 1$)
a mění stav jednoho skladu (vazba mění stav, multiplicita $* \rightarrow 1$).
Jeden uživatel může vytvořit libovolný počet pohybů (vazba vytváří, multiplicita $1 \rightarrow *$).
Pohyb nemění fyzickou lokaci součástky -- ta je pevně dána vazbou Lokace $\rightarrow$ Soucastka.

Vazba mezi Sklad a Soucastka je \textbf{kvalifikovaná} -- sklad přistupuje ke konkrétní
součástce prostřednictvím kvalifikátoru id: int.

\begin{figure}[H]
  \centering
  \makebox[\textwidth][c]{\includegraphics[width=1.2\textwidth]{02_class_diagram.png}}
  \caption{Class diagram}
  \label{fig:class_diagram}
\end{figure}

\newpage
\chapter{Dynamický model}

\section{Stavový diagram třídy SkladovyPohyb}

Třída SkladovyPohyb reprezentuje obecný skladový pohyb -- tedy jakoukoli změnu stavu zásob součástek
v systému. Každý pohyb prochází životním cyklem řízeným třemi metodami:
\begin{itemize}
  \item vytvoritPohyb() -- inicializuje nový pohyb: přiřadí ID, nastaví datum vzniku, typ pohybu a vazbu na součástku;
  \item overitData() -- ověří platnost zadaných dat: existenci součástky, vyplnění povinných polí,
        dostatečné množství na skladě (příjem/výdej) nebo platnost lokací (přesun);
  \item provestPohyb() -- fyzicky provede pohyb podle typu:
    \begin{itemize}
      \item příjem $\rightarrow$ soucastka.zvysMnozstvi(pocet)
      \item výdej $\rightarrow$ soucastka.snizMnozstvi(pocet)
      \item přesun $\rightarrow$ zmenLokaci()
    \end{itemize}
\end{itemize}
\begin{figure}[H]
  \centering
  \makebox[\textwidth][c]{\includegraphics[width=1.2\textwidth]{09_state_pohyb_skladu.png}}
  \caption{Stavový diagram třídy SkladovyPohyb}
  \label{fig:state_pohyb}
\end{figure}

\newpage
\section{Stavový diagram třídy Soucastka}

Diagram zachycuje, jak se mění stav součástky v~průběhu skladových operací.
Součástka je přidána do systému a okamžitě vstupuje do stavu Dostupna s~nulovým množstvím.
Stav Dostupna je hlavní pracovní stav součástky -- může mít mnozstvi = 0 (čeká na naskladnění)
i nenulové množství, přičemž část může být blokována atributem rezervovaneMnozstvi.
Metoda overDostupnost() ověří, zda je volné množství (mnozstvi - rezervovaneMnozstvi > 0).
Stav součástky je řízen těmito metodami:
\begin{itemize}
  \item zvysMnozstvi(pocet) -- zavolá se při naskladnění; zvýší atribut mnozstvi, součástka zůstává ve stavu Dostupna.
  \item rezervovat(pocet) -- zablokuje část kusů pro projekt; zvýší rezervovaneMnozstvi, stav se nemění.
  \item zrusitRezervaci(pocet) -- uvolní rezervaci; sníží rezervovaneMnozstvi, stav se nemění.
  \item snizMnozstvi(pocet) -- zavolá se při výdeji ze skladu (Vydana) nebo odpisu (Odepsana); oba jsou terminální stavy.
\end{itemize}
\begin{figure}[H]
  \centering
  \includegraphics[width=0.85\textwidth]{10_state_soucastka_projekt.png}
  \caption{Stavový diagram třídy Soucastka}
  \label{fig:state_soucastka_projekt}
\end{figure}

\newpage
\chapter{Model interakcí}

\section{Use Case diagram}

\textbf{Vedoucí skladu} má na starosti správu součástek a projektů a zadává skladové pohyby. \textbf{Skladník} fyzicky provádí příjem, výdej a přesun součástek a kontroluje stav zásob.

\begin{figure}[H]
  \centering
  \makebox[\textwidth][c]{\includegraphics[width=0.95\textwidth]{01_use_case.png}}
  \caption{Use Case diagram}
  \label{fig:use_case}
\end{figure}

\newpage
\section{Scénáře}

\subsection{Scénář: Příjem na sklad (PrijemPohyb)}
\begin{enumerate}
  \item Skladník zvolí typ pohybu příjem a zadá součástku, množství a dodavatele.
  \item Systém volá overitData() -- ověří existenci součástky a platnost cílové lokace.
  \item Pokud data OK $\rightarrow$ systém volá provestPohyb(), který interně volá Soucastka.zvysMnozstvi().
  \item Pohyb je potvrzen a zásoby jsou aktualizovány.
  \item \textbf{Alternativní tok:} součástka neexistuje nebo lokace je neplatná $\rightarrow$ systém vrátí chybovou hlášku a pohyb je zamítnut.
\end{enumerate}

\subsection{Scénář: Výdej ze skladu (VydejPohyb)}
\begin{enumerate}
  \item Skladník zvolí typ pohybu výdej a zadá součástku, množství a důvod výdeje.
  \item Systém volá overitData() -- ověří existenci součástky a zda mnozstvi $\geq$ požadovanému množství.
  \item Pokud data OK $\rightarrow$ systém volá provestPohyb(), který interně volá Soucastka.snizMnozstvi().
  \item Pohyb je potvrzen a zásoby jsou sníženy.
  \item \textbf{Alternativní tok:} nedostatečné množství na skladě $\rightarrow$ systém vrátí chybovou hlášku a pohyb je zamítnut.
\end{enumerate}

\subsection{Scénář: Přesun součástky (PresunutiPohyb)}
\begin{enumerate}
  \item Skladník zvolí typ pohybu přesun a zadá součástku, množství, zdrojovou a cílovou lokaci.
  \item Systém volá overitData() -- ověří existenci součástky, dostatečné množství ve zdrojové lokaci a platnost cílové lokace.
  \item Pokud data OK $\rightarrow$ systém volá provestPohyb() a poté PresunutiPohyb.zmenLokaci().
  \item Pohyb je potvrzen; součástka je evidována v~cílové lokaci.
  \item \textbf{Alternativní tok:} zdrojová nebo cílová lokace je neplatná $\rightarrow$ systém vrátí chybovou hlášku a pohyb je zamítnut.
\end{enumerate}

\subsection{Scénář: Přiřazení součástky k~projektu}
\begin{enumerate}
  \item Vedoucí skladu vybere projekt a zadá součástku s~množstvím.
  \item Systém volá Soucastka.overDostupnost() -- ověří, zda mnozstvi $-$ rezervovaneMnozstvi $> 0$.
  \item Pokud je dostatek volných kusů $\rightarrow$ systém volá Projekt.priradSoucastku(), které interně zavolá Soucastka.rezervovat().
  \item Součástka je přiřazena k~projektu a rezervované množství je aktualizováno.
  \item \textbf{Alternativní tok:} nedostatečný počet volných kusů $\rightarrow$ systém vrátí chybovou hlášku a přiřazení je zamítnuto.
\end{enumerate}

\subsection{Scénář: Správa součástek}
\begin{enumerate}
  \item Vedoucí skladu přidá novou součástku nebo upraví existující.
  \item Systém ověří validitu zadaných dat (název, počáteční množství, lokace).
  \item Pokud data OK $\rightarrow$ součástka je uložena; zásoby jsou aktualizovány.
  \item \textbf{Alternativní tok:} neplatná data $\rightarrow$ systém vrátí chybovou hlášku, součástka není uložena.
\end{enumerate}

\section{Sekvenční diagram -- Vytvoření pohybu}

Uživatel spustí vytvoření pohybu voláním vytvoritPohyb() na :SkladovyPohyb. Ten nejprve ověří platnost dat -- pokud data nesedí, vrátí zamítnutí. Při úspěšném ověření se zavolá provestPohyb(), které podle typu pohybu volá zvysMnozstvi(), snizMnozstvi() nebo overDostupnost() a zmenLokaci(). Uživatel dostane potvrzení po úspěšném dokončení.

\begin{figure}[H]
  \centering
  \makebox[\textwidth][c]{\includegraphics[width=1.2\textwidth]{13_sequence_pohyb.png}}
  \caption{Sekvenční diagram -- Vytvoření pohybu}
  \label{fig:sequence_pohyb}
\end{figure}

\newpage
\section{Diagram aktivit -- Vytvoření pohybu}

Diagram zachycuje průběh vytvoření skladového pohybu v~systému.
Pohyb prochází třemi fázemi řízenými metodami třídy SkladovyPohyb:
\begin{itemize}
  \item vytvoritPohyb() -- uživatel zadá typ pohybu, součástku a množství; metoda inicializuje pohyb, přiřadí ID a datum vzniku.
  \item overitData() -- ověří platnost zadaných dat: existenci součástky, dostatečné množství na skladě a platnost lokací u~přesunu.
  \item provestPohyb() -- fyzicky provede pohyb podle typu: příjem zavolá zvysMnozstvi(pocet), výdej snizMnozstvi(pocet), přesun zmenLokaci(). Pokud data nejsou platná, pohyb je zamítnut a uložena chyba.
\end{itemize}

\begin{figure}[H]
  \centering
  \makebox[\textwidth][c]{\includegraphics[width=1.2\textwidth]{12_activity_wms.png}}
  \caption{Diagram aktivit -- Vytvoření pohybu}
  \label{fig:activity_wms}
\end{figure}

\newpage
\chapter{Závěr}

Projekt se zaměřil na návrh informačního systému pro správu skladu elektronických součástek pomocí UML diagramů.
Výsledkem je class diagram s generalizací, kvalifikovanou vazbou, agregací a atributem spojení, dva stavové diagramy a model interakcí zahrnující Use Case, scénáře, sekvenční diagram a diagram aktivit.
Jednotlivé modely na sebe navazují a popisují stejný systém z různých pohledů.
Projekt pokrývá základní část systému -- evidenci součástek, pohybů a projektů -- a může být rozšířen o správu objednávek nebo inventurní modul.
Umělá inteligence (Claude) pomohla s formátováním a sazbou dokumentu v LaTeXu, urychlila tvorbu PlantUML diagramů a pomohla upřesnit finální podobu dokumentu.


\appendix
\chapter{Zdrojové kódy PlantUML diagramů}

\section{Use Case diagram}
\lstinputlisting[language={}, basicstyle=\ttfamily\scriptsize, breaklines=true, frame=single, lineskip=-1pt]{../diagrams/01_use_case.puml}

\newpage
\section{Class diagram}
\lstinputlisting[language={}, basicstyle=\ttfamily\scriptsize, breaklines=true, frame=single, lineskip=-1pt]{../diagrams/02_class_diagram.puml}

\newpage
\section{Stavový diagram – SkladovyPohyb}
\lstinputlisting[language={}, basicstyle=\ttfamily\scriptsize, breaklines=true, frame=single, lineskip=-1pt]{../diagrams/09_state_pohyb_skladu.puml}

\newpage
\section{Stavový diagram – Soucastka}
\lstinputlisting[language={}, basicstyle=\ttfamily\scriptsize, breaklines=true, frame=single, lineskip=-1pt]{../diagrams/10_state_soucastka_projekt.puml}

\newpage
\section{Sekvenční diagram – Vytvoření pohybu}
\lstinputlisting[language={}, basicstyle=\ttfamily\scriptsize, breaklines=true, frame=single, lineskip=-1pt]{../diagrams/13_sequence_pohyb.puml}

\newpage
\section{Diagram aktivit – Vytvoření pohybu}
\lstinputlisting[language={}, basicstyle=\ttfamily\scriptsize, breaklines=true, frame=single, lineskip=-1pt]{../diagrams/12_activity_wms.puml}
